% ============================================================
\section{Protocol and Definitions}
\label{sec:method}
% ============================================================

\textbf{Moirai for forecasting.}
Moirai~\cite{woo2024moirai} predicts a probabilistic distribution
over a target horizon given context.  For ETT forecasting we report
MSE on the median of 20 forecast samples at $h{=}96$ and $h{=}192$
(Appendix~\ref{app:moirai_median}); the IoT anomaly score of the
negative control is defined in Appendix~\ref{app:cnn_detail}.

\textbf{Training details.}
AdamW ($\eta = 10^{-4}$, weight decay $10^{-2}$), batch size 32,
early stopping with patience 3, gradient clipping at 1.0.  ETT
fine-tuning runs 20 epochs at $n{=}500$ (10 seeds) and is swept
across $n{\in}\{500,1{\text{k}},2{\text{k}},5{\text{k}},10{\text{k}}\}$.
\textbf{Reproducing any Moirai fine-tune requires the \texttt{uni2ts}
\texttt{PackedStdScaler} patch of Appendix~\ref{app:scaler_patch}}, without which
training can appear to run while silently not propagating gradients.

% Review M8: the appendix reports mean-pooled against token-flattened CKA, per-layer profiles and a
% random-init floor, but the body never said which of them its CKA column IS. Named here once, and with
% no new numerals beyond the window count: the flattened and per-layer values stay in the appendices
% this points at, because restating them here costs printed lines and a second set of chk() sites for
% numbers already covered there. The window count is deliberately NOT given: it is 300 on the Moirai
% arms and 200 on the non-Moirai ones, so a single number here would be wrong for a third of the matrix,
% and the per-arm counts are in the appendices with the protocols they belong to.
\textbf{Representation diagnostics, and which representation is primary.}  Two
metrics: linear CKA~\cite{kornblith2019cka} between pre-trained and fine-tuned
encoder representations, and weight drift ($\ell_2$ distance from pre-trained
parameters).  \textbf{Every CKA in the body is of the final encoder output,
mean-pooled over tokens}, on a fixed prefix of the selection split---the
comparison a practitioner holding two checkpoints can actually make, and the one
variant that needs no choice of layer or token position.  The choice is not free:
computed token-flattened instead, the same encoders read uniformly high and the
per-layer profile compresses too far to order.  That variant, the per-layer
profiles and a random-initialisation floor are all reported
(Appendices~\ref{app:layerunfreeze},~\ref{app:randominit}).

% Review S13/S19: name the screen precisely once, here, and document "gate" as its short form rather
% than renaming 276 sites. "Capability worth preserving" is defined in this paragraph's first sentence
% and used as the shorthand from here on, which is what S19 asked for.
\textbf{Baseline-relative pre-trained advantage.}  Nothing can be said about a
checkpoint having a \emph{capability worth preserving} except relative to a
baseline, so the baseline belongs in the quantity.  For a baseline
$b$, $\Vb{b} = 1-\text{MSE}_\text{ZS}/\text{MSE}_b$ is positive when the
pre-trained model beats $b$, and thresholding it at $0.20$---a prespecified
operating point, not an established boundary---is the \emph{baseline-relative
capability screen}, for which we write \emph{gate} in tables, figures and code.
\textbf{$\Vb{b}$ is a property of the pair, not of the checkpoint}: every
count we report is a count against a named $b$, and
\S\ref{sec:forecasting:gate} reports a ladder of eight.
% Review M3 asked for the consequence of baseline-relativity to be said out loud rather than left to be
% inferred from the property-of-the-pair sentence above it. Two lines, and they are what make the
% eight-rung ladder a necessary part of the design rather than a robustness check.
\emph{Capability worth preserving} is therefore not a property a checkpoint
\emph{has}: a claim of that form would be a claim against every baseline a
practitioner might deploy instead, and the ladder is the closest this design comes
to making one.  Our primary $b$ is one
ridge map $\R^{96 \times F} \to \R^{h \times F}$ ($\lambda{=}10^{-4}$) fit by
least squares on the \emph{training} windows and applied unchanged to the
evaluation windows, on the train-normalised scale the runs report (up to 300
windows).  Both terms see the same targets, and the pre-trained model sees
$96{+}h$ steps of context against the baseline's $96$, which if anything makes
$\Vb{b}$ generous to the pre-trained model.

\textbf{Three splits, and which quantity each one scores.} The \emph{training}
split fits both the fine-tuning updates and every baseline; the \emph{selection}
split chooses checkpoints and early stopping \textbf{and is where we score
$\Vb{b}$}; the \emph{held-out} split is disjoint from both and is where every
outcome is scored.  So a cell's admission cannot borrow the windows that later
judge it.  We also report the retrospective variant---$\Vb{b}$ on the held-out
windows themselves---which admits two fewer cells and changes no conclusion
(Appendices~\ref{app:gatecorrection},~\ref{app:corrections}).

\textbf{The estimand, named once.} For a cell with $S$ paired seeds,
\begin{equation}
\label{eq:denc}
\denc \;=\; \frac{1}{S}\sum_{s=1}^{S} 100 \cdot
\frac{\mathcal{L}_\text{test}(\text{B}_s) - \mathcal{L}_\text{test}(\text{D}_s)}
     {\mathcal{L}_\text{test}(\text{ZS})},
\end{equation}
the difference in forgetting between full fine-tuning~(B) and the frozen
encoder~(D), positive when freezing does better.  Both arms share the zero-shot
denominator, so $\denc$ is the raw contrast
$\mathcal{L}_\text{test}(\text{B})-\mathcal{L}_\text{test}(\text{D})$ rescaled per
cell to make different loss scales comparable, with no sign or ordering changed.
It is the estimand because it is the contrast this design \emph{manipulates}: B
and D differ in exactly one thing, whether the encoder's weights may move.
\textbf{Representation drift is not manipulated.}  CKA is an observational
\emph{post-treatment} measurement---taken after the intervention, on the object
the intervention changed---so it enters only as a candidate \emph{predictor} of
$\denc$.  That asymmetry is why a CKA-to-outcome association would not license a
CKA-based decision rule even had we found one.

% Round 7: the registered clauses are UNCHANGED -- moving a pre-registered definition after seeing
% outcomes is the thing the registration exists to prevent -- but a unanimity conjunction is a count
% of agreeing seeds, not an effect with an uncertainty, and the reviewer is right that a headline
% cannot rest on it. So the interval form is stated beside it and made primary, and both are reported.
% The "unanimity because an exact test cannot reach p<0.05 on 28 of 31 cells" justification is deleted
% from here rather than reworded: app:pairedinference argues it at length under its own heading, and
% those lines are what pays for the sentences added here.
\textbf{Definition (degradation of a demonstrated capability), and its interval
form.} A cell counts as degradation when all three clauses hold: (i)~$\Vb{b} \ge
0.20$ on selection windows, so there is a measured pre-trained advantage over
$b$; (ii)~full fine-tuning ends worse than not fine-tuning at all,
forg.$_\text{B} > 0$; and (iii)~the frozen encoder ends better,
forg.$_\text{D} < 0$, localising the loss to encoder weight updates rather than
to the data or the epoch budget. \textbf{Clauses (ii) and (iii) are reported two
ways}: as pre-registered, each holds in \emph{every} seed; and---primary from
this revision---as a paired interval, the $95\%$ CI on
$\text{B}{-}\text{ZS}$ lying entirely above zero and the CI on
$\text{D}{-}\text{ZS}$ entirely below, propagating the zero-shot reference's own
error and adjusted for multiplicity (Appendix~\ref{app:pairedinference}).
Intervals are primary because unanimity measures agreement among seeds rather
than the size of an effect or our confidence in it; the registered form stays
because retiring a criterion after seeing its outcome is not a choice we get to
make. \textbf{Both return the same count}, so nothing in this paper turns on
which is read. The definition is fixed for the whole paper.

\textbf{Conditions.} \textbf{B}~NLL-only fine-tune (baseline: how much does
unconstrained fine-tuning drift and forget?); \textbf{C}~NLL$+$SupCon
(Appendix~\ref{app:full_tables});
\textbf{D}~frozen encoder (control: separates encoder updates from head
training and overfitting); \textbf{E}~LoRA ($r{=}8$); \textbf{F}~L2-SP and
\textbf{G}~EWC. Conditions~D and~B are the comparison every claim in this
paper turns on.  \textbf{We name conditions rather than letter them from here on}
---\emph{full fine-tuning} (B), \emph{frozen encoder} (D), \emph{strict frozen
encoder} (H), \emph{LoRA} (E)---and keep the letters for the tables, the clauses
of the definition above and Equation~\ref{eq:denc}, so that no result in the body
has to be read through a lookup.

\textbf{What condition~D does and does not hold fixed.} D freezes the encoder's
own weights but leaves the input projection and mask encoding trainable, so it
controls \emph{encoder weight updates} and not representations: its measured CKA
is near but never exactly $1$ on Moirai ($0.969$--$0.9999$) and falls as low as
$0.38$ on TimesFM, where the encoder stack's weight drift is exactly zero.
\textbf{H}~(strict freeze) also freezes those two modules, is the only condition
with CKA exactly $1.000$, and reverses $\denc$'s sign on 4 of the 16 cells we ran
it on---every Moirai/ETTm2 cell.  Figure~\ref{fig:freezeboundary} draws which
modules move under each (Appendix~\ref{app:strictfreeze}).

\textbf{Dispersion.} Every ${\pm}$ in the body and in every generated table is
the standard error of the mean over seeds; where an appendix table or passage
uses a standard deviation instead it says so, and says which, where it does.

% The two lessons of this round, boxed because they are the parts of this paper most likely to be
% useful to someone who disagrees with its conclusion: both are choices we got wrong first, and both
% move results against us. The prose that used to state (1) in S7 and (2) in S4 is now a pointer, so
% the box costs its frame and not its content. No package added -- fbox+minipage only; the preamble is
% at its limit and a new float or tcolorbox risks a reflow the page budget cannot absorb.
\begin{center}
\fbox{\begin{minipage}{0.95\linewidth}
\small
\textbf{Two measurement choices that decide this result.}
\emph{(1)~Score the frozen-encoder comparison on windows you did not select on.}
Scoring $\denc$ on the selection windows instead reverses its sign on 4 of our 31
cells, and all four reverse toward flattering the frozen encoder
(\S\ref{sec:analysis:heldout}).
\emph{(2)~Fit the linear baseline.}  An earlier version of this work reached the
opposite conclusion from a baseline that was an unfitted per-window trend
extrapolation---a denominator a constant predictor beats---and correcting the
estimator moves 17 of Moirai's 21 cells across the threshold, 16 of them to
gate-fail (Appendix~\ref{app:gatecorrection}).
\end{minipage}}
\end{center}
